% ============================================================
% chapter2-Search_Pipeline.tex
% ============================================================
\section{La Pipeline di Elaborazione della Query}

La Search Pipeline costituisce il nucleo operativo dell'assistente ed è
orchestrata interamente all'interno del modulo \path{search_pipeline.py}.
Alla ricezione di una richiesta in linguaggio naturale, il sistema attiva
una sequenza strutturata di elaborazioni che sfrutta in modo sistematico
la concorrenza asincrona offerta da \path{asyncio}, combinando interrogazioni
parallele a database interni ed API esterne per produrre risultati
classificati, arricchiti e contestualizzati. L'intera pipeline si articola
in undici fasi distinte, ciascuna delle quali alimenta le successive con
strutture dati progressivamente più raffinate.

\vspace{0.4cm}
\noindent
\textbf{1. Recupero Anticipato della Memoria (Phantom Fetch)}\\
Non appena il testo grezzo della richiesta raggiunge la pipeline, viene
immediatamente avviata in background un'operazione di interrogazione su
\textbf{Qdrant}, il database vettoriale del sistema, tramite il metodo
\path{retrieve_context}. Questa operazione, denominata \emph{Phantom Fetch},
non attende il completamento delle fasi di analisi semantica: agisce sul
testo non ancora analizzato per recuperare, con il massimo anticipo
temporale, le inserzioni storiche già validate o consigliate in sessioni
precedenti su query semanticamente affini. Il principio sottostante è che
prodotti già incontrati e giudicati positivamente rappresentano una
baseline informativa di alta qualità, e il loro recupero anticipato
consente di parallelizzare il costo di latenza con le fasi computazionalmente
più intensive che seguono.

\vspace{0.4cm}
\noindent
\textbf{2. Recupero dello Storico Conversazionale a Breve Termine}\\
Per consentire all'assistente di mantenere la coerenza su conversazioni
multi-turno, nelle quali l'utente può affinare progressivamente i propri
criteri, la pipeline recupera la cronologia recente della sessione attiva tramite \path{redis_client.get_user_queries}.
Vengono estratte le ultime tre query (\path{history[:3]}), costruendo uno
\textit{snapshot} contestuale volutamente leggero: la Search Pipeline
necessita di un riferimento minimo per risolvere anafore e riferimenti
impliciti, senza sovraccaricare la \textit{context window} dell'LLM
impiegato nella fase di parsing. La gestione della memoria a lungo
termine, incluse le eventuali logiche di \textit{summarization},
è deliberatamente demandata all'Orchestratore ReAct, circoscrivendo
il compito della pipeline al solo \textit{intent-matching} istantaneo
sulla sessione corrente.

\vspace{0.4cm}
\noindent
\textbf{3. Analisi Logica della Query (NLP Parsing)}\\
Il testo grezzo e lo storico conversazionale vengono consegnati al servizio
\path{parse_query_service}, responsabile della trasformazione del linguaggio
naturale in strutture rigide e interrogabili dai sistemi a valle. Il processo
si articola in due livelli complementari. 
Il primo livello è un'estrazione rapida basata su \textbf{spaCy}: il modello
individua in modo deterministico elementi lessicali ad alta rilevanza, come
soglie numeriche di budget (\emph{"sotto i 200 euro"} genera il vincolo
\path{max_price=200}) e indicatori di condizione (\emph{"rovinato"} forza
il filtro \path{condition="usato"}). Questo stadio garantisce latenza
minima per i casi più frequenti e ben strutturati.
Il secondo livello delega a un LLM l'interpretazione delle
ambiguità semantiche non risolvibili con regole lessicali. Il modello
restituisce un oggetto strutturato che include il marchio, le specifiche
tecniche e il campo \path{intent_product_type}: un classificatore binario 
che distingue se l'utente intende acquistare un dispositivo principale (\texttt{Main}) 
o un accessorio (\texttt{Accessory}). Questa distinzione governa le fasi di filtraggio successive.
In parallelo, il modulo \path{_build_ebay_query} effettua una traduzione
sistematica delle negazioni in linguaggio naturale italiano, ad esempio espressioni
come \emph{"senza cover"}, \emph{"no usati"}, \emph{"nessun graffio"},
nella sintassi di esclusione supportata dalle API di ricerca di eBay,
prependendovi il segno aritmetico meno (\texttt{-}). Questa sostituzione
è fondamentale per evitare che termini negativi del linguaggio naturale
vengano interpretati come keyword di inclusione dai motori di ricerca esterni.

\vspace{0.4cm}
\noindent
\textbf{4. Ricerca Parallela API eBay ed Espansione Semantica Locale (RAG)}\\
La pipeline lancia simultaneamente, tramite \path{asyncio.gather}, due
flussi di recupero indipendenti che operano su sorgenti dati eterogenee.
Da un lato, \path{search_items} interroga le \textbf{Browse API di eBay},
richiedendo un sottoinsieme ristretto delle inserzioni più
rilevanti secondo i criteri costruiti nella fase precedente (massimo 15 risultati). La scelta
di limitare il volume dei risultati è deliberata: un insieme pre-selezionato 
consente elaborazioni successive più rapide rispetto a un ampio dataset rumoroso.
Dall'altro lato, si attiva il modulo \path{_do_rag}, che implementa un approccio \textbf{Retrieval-Augmented Generation} sul corpus documentale interno. La query viene prima espansa semanticamente tramite \path{expand_query}, che identifica sinonimi e varianti terminologiche per massimizzare la copertura del recupero. L'interrogazione si biforca quindi su due collezioni Qdrant distinte: \path{product_docs}, contenente schede tecniche e documentazione di prodotto, e \path{seller_docs}, che raccoglie guide operative e casistiche di problemi ricorrenti segnalati dai venditori. L'integrazione di queste due sorgenti consente all'assistente di contestualizzare le inserzioni eBay con informazioni qualitative non presenti nelle schede commerciali.

\vspace{0.4cm}
\noindent
\textbf{5. Filtro Anti-Accessori (Candidate Pool Fallback)}\\
Un problema strutturale delle ricerche su marketplace generalisti è la prominenza degli accessori nei risultati per query relative a dispositivi principali: interrogando le API di eBay con un termine sintetico come \emph{"iPhone 13"}, la risposta di default può essere composta interamente da cover, pellicole protettive e caricatori.
Sebbene la pipeline disponga di un algoritmo di Reranking vettoriale a valle, affidarvisi per correggere questo degrado sarebbe architetturalmente scorretto. Il problema non riguarda l'ordinamento, bensì la \textbf{Candidate Generation}: un reranker, per quanto sofisticato, non può promuovere in prima posizione un dispositivo che non sia presente nel pool di candidati. Poiché il pool è deliberatamente limitato a un massimo di quindici inserzioni, è sufficiente che il corpus estratto contenga accidentalmente quattordici accessori per rendere il reranker incapace di restituire risultati utili.
Il filtro interviene pertanto a monte, analizzando i primi cinque risultati eBay: se almeno tre recano etichette categoriali riconducibili ad Accessori, Ricambi o Custodie (\emph{accessory dominance}) e l'intento classificato in Fase 3 è \texttt{Main}, l'intero sottoinsieme viene scartato. Il sistema esegue quindi il crawling degli \emph{Aspect Distributions} della prima pagina per recuperare dinamicamente il \path{categoryId} corrispondente alla sezione Elettronica pertinente, e invia una seconda richiesta alle API di eBay vincolata a quella categoria, garantendo che il pool per le fasi successive sia composto esclusivamente da dispositivi attesi.

\vspace{0.4cm}
\noindent
\textbf{6. Fusione Ibrida dei Risultati (Reciprocal Rank Fusion)}\\
I risultati provenienti dalla ricerca live su eBay vengono fusi con
il corpus storico recuperato durante il Phantom Fetch all'interno
del modulo \path{_merge_hybrid_results}. Prima della
fusione vera e propria, ogni inserzione storica viene sottoposta a
una verifica di conformità rispetto ai vincoli di budget attuali:
gli elementi che non soddisfano i limiti di prezzo correnti vengono
scartati, garantendo che il corpus storico non introduca risultati
obsoleti o non pertinenti alla sessione in corso.
Su questa base, il sistema applica un meccanismo di rinforzo
(\path{_hybrid_bonus}): se un prodotto risulta presente \emph{sia} nel
corpus storico \emph{sia} tra i risultati live di eBay, il suo score RRF
viene incrementato di \(+0.15\) punti. La motivazione è diretta: un
prodotto che compare spontaneamente in due sorgenti indipendenti,
costruite in momenti e con logiche diverse, ha maggiori probabilità
di essere genuinamente rilevante rispetto a uno che appare in una sola di esse.

\vspace{0.4cm}
\noindent
\textbf{7. Architettura del Sistema di Ranking (Reranker e Cross-Encoder)}\\
Questa fase concentra le operazioni di riordinamento semantico
più computazionalmente intensive dell'intera pipeline, affidandosi
a un'architettura di ranking a due stadi che combina feature
ingegnerizzate manualmente con un modello neurale di tipo Cross-Encoder.
L'implementazione attuale è il risultato di un percorso evolutivo
di tre stadi. Un primo approccio basato su \textbf{LightGBM},
addestrato sull'intero dataset SIGIR Coveo E-commerce,
ha prodotto performance insoddisfacenti a causa della scarsa qualità
intrinseca dei segnali di click. Questa evidenza ha orientato lo sviluppo verso l'architettura attualmente in produzione.
Il processo ha inizio con il pre-calcolo del \path{trust_score} di
ciascun venditore tramite \path{_prefetch_top_sellers_feedback}: i
feedback vengono scaricati in parallelo e sottoposti ad
algoritmi di \textit{sentiment analysis}, producendo un punteggio
reputazionale aggregato. Ottenuti questi valori, la pipeline invoca il modulo di
\textbf{Vector Reranking} (\path{rerank_products}), che opera
attraverso due stadi in cascata:
Il primo è uno \textbf{stadio Base-Score} a ponderazione pesata,
che valuta ogni inserzione su un sistema di oltre quindici feature
ingegnerizzate manualmente, tra cui figurano la similarità
semantica (\path{semantic_sim}), la similarità lessicale
(\path{lexical_sim}), il punteggio di fiducia (\path{trust_score}), la distanza normalizzata dal prezzo mediano
(\path{price_z}), il boost derivante dal corpus RAG
(\path{rag_product_boost}) e il punteggio di accessorietà
(\path{accessory_score}). Il secondo è uno \textbf{stadio
Cross-Encoder vettoriale}, applicato ai soli dieci risultati
provvisoriamente in testa per ammortizzare il costo computazionale:
il modello \path{ms-marco-MiniLM-L-6-v2} genera un'inferenza
profonda (\emph{cross-score}) tra il testo della query originale
e i documenti \path{product_docs}. I due score vengono infine combinati mediante un blending
pesato che assegna il \textbf{70\%} del peso al cross-score
vettoriale e il restante \textbf{30\%} al base-score,
privilegiando la componente semantica senza eliminare il contributo delle feature deterministiche.
Nel medesimo intervallo temporale operano in background tre micro-moduli: \path{get_return_policies} per identificare l'idoneità ai resi della categoria dominante, \path{update_user_profile} per aggiornare le affinità del cliente su Postgres in base alla chiamata odierna e \path{extract_attributes_batch} che esegue la \textbf{Named Entity Recognition (NER)} estraendo attributi base dai titoli e pulendoli.

\vspace{0.4cm}
\noindent
\textbf{8. Persistenza Transazionale Preliminare}\\
Per isolare la latenza fisiologica indotta dalle operazioni sul
database relazionale dal percorso critico della classificazione
finale, la memorizzazione permanente viene avviata prima ancora
che la graduatoria euristica sia conclusa.
Il modulo \path{_prepare_and_persist_items} opera in un \textit{thread}
isolato, effettuando i salvataggi necessari sulla tabella Postgres.
\textbf{Redis} agisce a monte come layer di deduplicazione: tramite
la chiave \path{known_ebay_ids}, il sistema intercetta ogni
operazione di scrittura e blocca gli identificatori eBay già
presenti nell'archivio, inviando in \texttt{INSERT} unicamente le
inserzioni inedite, scongiurando \textit{bottleneck} e collisioni.

\vspace{0.4cm}
\noindent
\textbf{9. Classificazione Euristica Finale (Score Composito)}\\
Ogni inserzione sopravvissuta ai controlli precedenti riceve in
questa fase la sua calibrazione definitiva, calcolata da
\path{_apply_final_ranking}. All'ordinamento vettoriale prodotto
dal reranker nella Fase 7 vengono sovrapposti dei bonus deterministici 
cuciti sulle preferenze in memoria dell'utente che invia la query.
Il primo modificatore è il \textbf{Brand Affinity Bonus}: se il
database utente documenta una predilezione storica per un
determinato produttore, le inserzioni appartenenti a tale brand
ricevono un incremento di \(+0.15\) punti. Il secondo è costituito
dal \textbf{Condition Match Bonus}: se lo stato dell'oggetto 
collima con l'abitudine d'acquisto registrata dell'utente (ad es. predilezione per il \emph{nuovo} 
o il \emph{ricondizionato}), scatta un incentivo calcolato di \(+0.08\). In ultimo interviene 
un'\textbf{euristica sul value-for-money}: laddove il prezzo stimato si posizioni all'interno
o al di sotto delle proiezioni di budget tollerate dall'utente 
(\path{contextual_budgets}), viene erogato un premio addizionale di \(+0.10\) punti, 
premiando i prodotti che assecondino un superbo rapporto qualità/spesa per quel profilo.

\vspace{0.4cm}
\noindent
\textbf{10. Valutazione Qualitativa e Comparativa tramite LLM Discriminator}\\
Gli otto risultati in cima alla graduatoria vengono inviati al
modulo \path{_compute_llm_value_scores} per ricevere una
validazione in linguaggio naturale. A differenza di
un utilizzo esplorativo, il discriminatore opera su un \textit{market scope}
ristretto e dedotto esclusivamente dalle misurazioni della SERP eBay corrente:
riceve il prezzo mediano calcolato a runtime, l'escursione finanziaria e le distribuzioni delle condizioni.
Su questa base rigorosa, il modello formula per ciascun item uno \path{score} fluttuante da \(-1.0\) a \(1.0\) 
ed un giudizio in compendio (es: \emph{"💡 Eccellente condizione hardware del touch..."}), inibendo allucinazioni esterne.
In successione una invocazione analoga (\path{explain_results})
analizza unicamente i primi \textbf{tre risultati primatisti}: ne distilla una valutazione 
mirata stilando elenchi dettagliati di Pro e Contro (\emph{Pros \& Cons}). 
A conclusione organica della fase, vengono innestate le referenze dirette alle stringhe di policy (RAG Attach) scoperte in Fase 7.

\vspace{0.4cm}
\noindent
\textbf{11. Metriche IR e Indicizzazione Finale su Qdrant}\\
A conclusione dell'elaborazione, la pipeline formalizza in memoria
le metriche di qualità del recupero valutate secondo gli stringenti standard operativi
dall'\textbf{Information Retrieval}: in questa istanza sono approntate misurazioni formali su
\textbf{Precision@5}, \textbf{Recall@10} e \textbf{nDCG@10} (Normalized Discounted Cumulative Gain).
Tali metriche, esposte nel backend, saranno serbate per le diagnosi offline sulle performance dell'IR engine.
In chiusura del canale, i quindici elementi emersi vengono inseriti in pretesa allo spooler per l'indicizzazione (\path{index_search_items}).
Trascritti vettorialmente ed impacchettati su \textbf{Qdrant}, questi ultimi vanno a colmare stabilmente 
la riserva locale per le esecuzioni Phantom Fetch successive: la pipeline completa la propria esecuzione 
lasciando un'eredità circolare e resiliente valida a supportare le latenze del domani.

\vspace{1cm}
% ============================================================
% LTR EXPERIMENT
% ============================================================
\section{Approfondimento di Ricerca: Learning to Rank Personalizzato (LTR)}

Come richiamato in precedenza nell'evoluzione della pipeline (Fase 7), l'architettura di classificazione 
ha attraversato una fase di sviluppo intensivo mirata all'addestramento di un modello 
supervisionato nativo. Benché tale approccio sia stato scalzato in produzione dall'attuale sistema ibrido a Reranker, le complesse logiche simulative 
e ingegneristiche gettate per sopperire alle limitazioni del dominio costituiscono un impianto architetturale
di assoluto rilievo che merita documentazione formale.

\subsection{Obiettivo Generale e Scelta dell'Infrastruttura}
L'obiettivo iniziale era addestrare un modello di Learning to Rank (LTR) capace di riordinare i risultati di ricerca in base al netto profilo comportamentale dell'utente. A dispetto delle API primordiali standardizzate, il sistema mirava a comprendere che un consumatore votato al risparmio anela a posizioni economiche, slegandosi dalle pertinenze nominali. 
Il problema è stato formulato come ranking supervisionato: un modello tarato su \textbf{Pointwise LTR} in cui ogni coppia \emph{(utente, prodotto)} riceveva un'etichetta di rilevanza da $0$ a $N-1$, con addestramento pilotato su framework \textbf{LightGBM} (\path{rank:ndcg}).

A differenza di un \textit{cross-encoder} puro isolato, l'L2R permetteva il vaglio simultaneo dell'intera mole di prodotti in parallelo e forzava le \textit{feature demografiche/utente} all'interno del peso di predizione finale.

\subsection{Il Problema Critico e L'Ecosystem Simulation}
La criticità strutturale di tutto l'impianto è subentrata a causa della mancanza cronica di dataset pubblici di \textit{clickstream} o di giudizi di rilevanza massivi per eBay. Eseguire l'acquisizione su \textit{dwell time} o log umani si è rivelato proibitivo sia economicamente sia in termini d'utenza attiva reclutabile.

Per sormontare l'ostacolo, abbiamo progettato ed ingegnerizzato l'estrazione di un \textbf{Dataset Sintetico} per simulazione algoritmica (\textit{Ecosystem Simulaton}): uno strato di Agenti LLM delegati all'impersonificazione matematica di precisi profili comportamentali stereotipati (le \textit{Persona}). Tali profili erano vincolati a rigidità psicometriche:
\begin{enumerate}
    \item \textbf{Parametri Numerici}: Sensibilità assoluta al prezzo (\path{price_sensitivity}), affidabilità (\path{quality_preference}), condizione ambita (\path{condition_affinity}), fedeltà aziendale (\path{brand_loyalty}).
    \item \textbf{Limiti Qualitativi}: Un \textit{System Prompt} ideato per condizionare ferocemente lo stampo inferenziale del modello in assenza di scostamenti creativi.
\end{enumerate}

Le 10 \textit{Persona} sono state plasmate per marcare le estremità delle varianze demografiche d'impiego: profilazioni come il \textit{Bargain Hunter} (totale esasperazione del risparmio), \textit{L'Auction Sniper} (vigliacco ricercatore dell'affare occasionale), il \textit{Luxory Buyer} ed il prudente \textit{Safety Focused} esasperato dai feedback.

\subsection{Diversità Architetturale e Normalizzazione Estesa}
Un fattore decisionale saliente ai fini dell'onesta integrità algoritmica è scaturito dall'impedire un imprinting univoco da \emph{Foundation Models}. Un singolo LLM su dieci prompt avrebbe inquinato le estrazioni riversandovi pattern ed allucinazioni monolitiche legate all'architettura costruttrice di base. L'esecutivo logico è stato diviso frammentandone la paternità intellettuale: le elaborazioni asincrone venivano distribuite su \textbf{Groq Llama 3.3/3.1} (architettura Meta), \textbf{Qwen} ed \textbf{DeepSeek} (Cinese ad elevate asimmetrie concettuali), \textbf{Mistral-Nemo NVIDIA} e un'istanza locale \textbf{Ollama} a \textit{temperature} elevatissima a garanzia di entropia per profilazioni irrazionali (\emph{Emotional Bidder}).

Le iterazioni in serie processavano le richieste calcolando il valore \textit{prezzo/mediana-SERP} e le sottoponevano alle intelligenze con vincoli di ordinamento inverso (da $N-1$ d'eccellenza al fallimento pari a \(0\)). Le discrepanze numeriche scaturite — al netto di sistemi di ripiego sui \textit{Silent Fallback} dell'inferenza — venivano normalizzate tramite scale lineari assolute o scale logaritmiche (\(\log(1+x)/\log(1+p_{99})\)) al fine specifico di smorzare derive esponenziali di \textit{seller_feedback} milionari. 

\subsection{Integrazione LightGBM e Conclusioni Sperimentali}
Una volta depurate, le estrazioni fornivano una spina dorsale di \textbf{19 Colonne} e migliaia di inferenze \textit{Batch-Ranking} incrociate (\texttt{training\_set\_v2.csv}). Ne misurammo i tassi di collisione ricorrenze ed asimmetrie ideali mediante \textbf{Cohen's Kappa Quadratico Pesato} validando un \(\kappa \in [0.4, 0.6]\) quale sintomo di eccellenza eterogenea, iniettando la mole addestrativa finalmente nello spool \textit{lambdarank}:
\begin{verbatim}
import lightgbm as lgb
import pandas as pd

df = pd.read_csv("training_set_v2.csv")
groups = df.groupby(["query", "persona_id"]).size().values

dataset = lgb.Dataset(
    df[feature_cols],
    label=df["relevance"],
    group=groups
)

model = lgb.train(
    {"objective": "lambdarank", "metric": "ndcg", "ndcg_eval_at": [5, 10]},
    dataset
)
\end{verbatim}

A discapito delle premesse simulatrici e dell'eccellente distribuzione analitica implementata, i test pratici di validazione hanno evidenziato criticità insanabili sul mercato nativo: la propensione innata di qualsivoglia LLM a compiacere descrizioni organiche a discapito dell'istinto brutale da consumatore (\textit{residual bias in cross-models}), accavallati con l'assenza insormontabile di un \textit{gold set} confermativo di base organica su vasta scala. Pur abbandonando il \textit{lightgbm} puro in favore d'un asset ibrido (vedi Fase 7), si attesta e convalida un massiccio impianto formativo nella scalabilità della simulazione e della sintesi multi-rete per i \textit{Cold-start domain} di futura applicazione.
